
\documentclass{article}
\usepackage{colortbl}
\usepackage{makecell}
\usepackage{multirow}
\usepackage{supertabular}

\begin{document}

\newcounter{utterance}

\centering \large Interaction Transcript for game `function\_detective', experiment `list\_passive\_labeled\_pairs', episode 93 with Qwen3.6{-}35B{-}A3B{-}FP8{-}without{-}reasoning.
\vspace{24pt}

{ \footnotesize  \setcounter{utterance}{1}
\setlength{\tabcolsep}{0pt}
\begin{supertabular}{c@{$\;$}|p{.15\linewidth}@{}p{.15\linewidth}p{.15\linewidth}p{.15\linewidth}p{.15\linewidth}p{.15\linewidth}}
    \# & \multicolumn{2}{c}{Player} && \multicolumn{2}{c}{Game Master} \\
    \hline

    \theutterance \stepcounter{utterance}  
    & & & \multicolumn{4}{p{0.6\linewidth}}{
        \cellcolor[rgb]{0.9,0.9,0.9}{
            \makecell[{{p{\linewidth}}}]{
                \texttt{\tiny{[P1$\langle$GM]}}
                \texttt{I have a hidden Python function for you to figure out.} \\
\\ 
\texttt{Your job is to infer what it does and then submit a correct Python implementation.} \\
\\ 
\texttt{Function signature:} \\
\texttt{(items: List{[}int{]}) {-}> int} \\
\\ 
\texttt{You have 15 turns total.} \\
\\ 
\texttt{I reveal candidate input{-}output pairs one at a time and tell you whether each one is correct.} \\
\\ 
\texttt{You must make your best SOLVE attempt no later than the second{-}to{-}last turn.} \\
\texttt{Do not wait until the very end.} \\
\\ 
\texttt{You can use these actions:} \\
\texttt{{-} JSON format described at the end of the prompt} \\
\texttt{{-} SOLVE: ```python ... ```} \\
\\ 
\texttt{Do not add import statements in your guess.} \\
\texttt{Anything you might need is already available, so just write the function body in hypothesis and return the answer.} \\
\\ 
\texttt{Preloaded examples:} \\
\texttt{1. (({[}5, {-}5, 5{]}, 50), True)} \\
\\ 
\texttt{Scoring:} \\
\texttt{{-} correctness is required for any non{-}zero score} \\
\texttt{{-} in interactive modes, solving with fewer information requests is better} \\
\\ 
\texttt{Input formatting reminders:} \\
\texttt{{-} strings must be quoted} \\
\texttt{{-} booleans must be True/False} \\
\texttt{{-} use commas between multiple values} \\
\texttt{{-} do not write any extra action tags in the same message} \\
\\ 
\texttt{Format example:} \\
\texttt{You write:} \\
\texttt{\{'current\_mode': <probing/confirming/don't know; choose whichever best describes your current state in the task>, 'input\_rationale': <brief description of why you chose this hypothesis>, 'hypothesis': <either the exact string "unknown" if you have no guess yet OR your current best guess as runnable Python matching the given signature, with a one{-}line docstring stating what it does. Never write placeholder code, if you cannot commit to a guess, use "unknown". Match the given signature.} \\
\texttt{e.g.} \\
\texttt{def f(x):} \\
\texttt{	"""Returns the absolute value of x."""} \\
\texttt{	return x if x >= 0 else {-}x} \\
\\ 
\texttt{Use "unknown" if you have no guess.>, 'confidence\_score': <score between 1{-}100. 100 means highest confidence in your hypothesis>\}} \\
\\ 
\texttt{Output only one valid json per turn.} \\
            }
        }
    }
     \\ \\

    \theutterance \stepcounter{utterance}  
    & \multicolumn{4}{p{0.6\linewidth}}{
        \cellcolor[rgb]{0.561,0.737,0.561}{
            \makecell[{{p{\linewidth}}}]{
                \texttt{\tiny{[P1$\rangle$GM]}}
                \texttt{\{'current\_mode': 'probing', 'input\_rationale': 'I have one example: ({[}5, {-}5, 5{]}, 50) {-}> True. This suggests the function might check if a target sum (50) can be achieved by some combination of the items, possibly with signs. Let me probe with a simpler case to understand the behavior better.', 'hypothesis': 'unknown', 'confidence\_score': 10\}} \\
            }
        }
    }
    & & \\ \\

    \theutterance \stepcounter{utterance}  
    & & & \multicolumn{4}{p{0.6\linewidth}}{
        \cellcolor[rgb]{0.9,0.9,0.9}{
            \makecell[{{p{\linewidth}}}]{
                \texttt{\tiny{[P1$\langle$GM]}}
                \texttt{Game Violation: ParseError: Response is not valid JSON.} \\
\texttt{Output ONLY:} \\
\texttt{One JSON with confidence\_score, current\_mode, input\_rationale, and hypothesis keys ...} \\
\texttt{OR} \\
\texttt{SOLVE: ```python ... ```} \\
            }
        }
    }
     \\ \\

    \theutterance \stepcounter{utterance}  
    & \multicolumn{4}{p{0.6\linewidth}}{
        \cellcolor[rgb]{0.561,0.737,0.561}{
            \makecell[{{p{\linewidth}}}]{
                \texttt{\tiny{[P1$\rangle$GM]}}
                \texttt{\{} \\
\texttt{  "current\_mode": "probing",} \\
\texttt{  "input\_rationale": "I have one example: ({[}5, {-}5, 5{]}, 50) {-}> True. This suggests the function might check if a target sum (50) can be achieved by some combination of the items, possibly with signs. Let me probe with a simpler case to understand the behavior better.",} \\
\texttt{  "hypothesis": "unknown",} \\
\texttt{  "confidence\_score": 10} \\
\texttt{\}} \\
            }
        }
    }
    & & \\ \\

    \theutterance \stepcounter{utterance}  
    & & & \multicolumn{4}{p{0.6\linewidth}}{
        \cellcolor[rgb]{0.9,0.9,0.9}{
            \makecell[{{p{\linewidth}}}]{
                \texttt{\tiny{[P1$\langle$GM]}}
                \texttt{OUTPUT: (({[}{]}, {-}1), False)} \\
            }
        }
    }
     \\ \\

    \theutterance \stepcounter{utterance}  
    & \multicolumn{4}{p{0.6\linewidth}}{
        \cellcolor[rgb]{0.561,0.737,0.561}{
            \makecell[{{p{\linewidth}}}]{
                \texttt{\tiny{[P1$\rangle$GM]}}
                \texttt{\{} \\
\texttt{  "current\_mode": "probing",} \\
\texttt{  "input\_rationale": "The function takes a list of integers and a target integer. It returns True if the target can be formed by summing a subset of the items, possibly with signs flipped? Or maybe just subset sum? Let's test with a simple positive case. The first example ({[}5, {-}5, 5{]}, 50) {-}> True is strange because 5+({-}5)+5=5, not 50. Maybe it's about concatenation? Or maybe the second argument is not a target sum but something else? Wait, let me re{-}read. The signature is (items: List{[}int{]}) {-}> int. But the examples show two arguments. Let me check the preloaded example again: (({[}5, {-}5, 5{]}, 50), True). This looks like the function takes two arguments. But the signature says one. Perhaps the signature is wrong or I'm misinterpreting. Let me try a simple subset sum test.",} \\
\texttt{  "hypothesis": "unknown",} \\
\texttt{  "confidence\_score": 10} \\
\texttt{\}} \\
            }
        }
    }
    & & \\ \\

    \theutterance \stepcounter{utterance}  
    & & & \multicolumn{4}{p{0.6\linewidth}}{
        \cellcolor[rgb]{0.9,0.9,0.9}{
            \makecell[{{p{\linewidth}}}]{
                \texttt{\tiny{[P1$\langle$GM]}}
                \texttt{OUTPUT: (({[}0{]}, 0), True)} \\
            }
        }
    }
     \\ \\

    \theutterance \stepcounter{utterance}  
    & \multicolumn{4}{p{0.6\linewidth}}{
        \cellcolor[rgb]{0.561,0.737,0.561}{
            \makecell[{{p{\linewidth}}}]{
                \texttt{\tiny{[P1$\rangle$GM]}}
                \texttt{\{} \\
\texttt{  "current\_mode": "probing",} \\
\texttt{  "input\_rationale": "The function seems to take two arguments: a list of integers and a target integer. It returns True if the target can be formed by summing a subset of the items. Let's verify this hypothesis with a case where the target cannot be formed. For example, {[}1, 2{]} with target 4 should be False since 1+2=3 < 4 and no subset sums to 4.",} \\
\texttt{  "hypothesis": "unknown",} \\
\texttt{  "confidence\_score": 20} \\
\texttt{\}} \\
            }
        }
    }
    & & \\ \\

    \theutterance \stepcounter{utterance}  
    & & & \multicolumn{4}{p{0.6\linewidth}}{
        \cellcolor[rgb]{0.9,0.9,0.9}{
            \makecell[{{p{\linewidth}}}]{
                \texttt{\tiny{[P1$\langle$GM]}}
                \texttt{OUTPUT: (({[}1, 1{]}, 3), False)} \\
            }
        }
    }
     \\ \\

    \theutterance \stepcounter{utterance}  
    & \multicolumn{4}{p{0.6\linewidth}}{
        \cellcolor[rgb]{0.561,0.737,0.561}{
            \makecell[{{p{\linewidth}}}]{
                \texttt{\tiny{[P1$\rangle$GM]}}
                \texttt{\{} \\
\texttt{  "current\_mode": "probing",} \\
\texttt{  "input\_rationale": "The function appears to check if a target sum can be formed by a subset of the given integers. Let's test with a case where the target is exactly the sum of all elements to confirm subset sum behavior. For example, {[}1, 2{]} with target 3 should be True since 1+2=3.",} \\
\texttt{  "hypothesis": "unknown",} \\
\texttt{  "confidence\_score": 30} \\
\texttt{\}} \\
            }
        }
    }
    & & \\ \\

    \theutterance \stepcounter{utterance}  
    & & & \multicolumn{4}{p{0.6\linewidth}}{
        \cellcolor[rgb]{0.9,0.9,0.9}{
            \makecell[{{p{\linewidth}}}]{
                \texttt{\tiny{[P1$\langle$GM]}}
                \texttt{OUTPUT: (({[}1, 2, 3{]}, 14), True)} \\
            }
        }
    }
     \\ \\

    \theutterance \stepcounter{utterance}  
    & \multicolumn{4}{p{0.6\linewidth}}{
        \cellcolor[rgb]{0.561,0.737,0.561}{
            \makecell[{{p{\linewidth}}}]{
                \texttt{\tiny{[P1$\rangle$GM]}}
                \texttt{\{} \\
\texttt{  "current\_mode": "probing",} \\
\texttt{  "input\_rationale": "The previous hypothesis (subset sum) is incorrect because {[}1, 2, 3{]} sums to 6, not 14, yet the result is True. This suggests the function might involve multiplication, concatenation, or some other operation. Let's try a case with a single element to see if it's about divisibility or something similar. For example, ({[}2{]}, 4) might be True if it checks if the target is a multiple of the element, or False if it's strict subset sum. Let's try ({[}2{]}, 4).",} \\
\texttt{  "hypothesis": "unknown",} \\
\texttt{  "confidence\_score": 30} \\
\texttt{\}} \\
            }
        }
    }
    & & \\ \\

    \theutterance \stepcounter{utterance}  
    & & & \multicolumn{4}{p{0.6\linewidth}}{
        \cellcolor[rgb]{0.9,0.9,0.9}{
            \makecell[{{p{\linewidth}}}]{
                \texttt{\tiny{[P1$\langle$GM]}}
                \texttt{OUTPUT: (({[}{-}1, 0, 1{]}, 3), False)} \\
            }
        }
    }
     \\ \\

    \theutterance \stepcounter{utterance}  
    & \multicolumn{4}{p{0.6\linewidth}}{
        \cellcolor[rgb]{0.561,0.737,0.561}{
            \makecell[{{p{\linewidth}}}]{
                \texttt{\tiny{[P1$\rangle$GM]}}
                \texttt{\{} \\
\texttt{  "current\_mode": "probing",} \\
\texttt{  "input\_rationale": "The function is not simple subset sum. Let's look at the examples again: ({[}5, {-}5, 5{]}, 50) {-}> True, ({[}{]}, {-}1) {-}> False, ({[}0{]}, 0) {-}> True, ({[}1, 1{]}, 3) {-}> False, ({[}1, 2, 3{]}, 14) {-}> True, ({[}{-}1, 0, 1{]}, 3) {-}> False. The case ({[}1, 2, 3{]}, 14) {-}> True is very strange for subset sum. Maybe it's about forming the target using arithmetic operations on the numbers? Or maybe it's about the product? 1*2*3=6 != 14. Maybe it's about concatenation? No, they are ints. Let's try a case with a single large number to see if it's about divisibility or equality. Let's try ({[}14{]}, 14).",} \\
\texttt{  "hypothesis": "unknown",} \\
\texttt{  "confidence\_score": 30} \\
\texttt{\}} \\
            }
        }
    }
    & & \\ \\

    \theutterance \stepcounter{utterance}  
    & & & \multicolumn{4}{p{0.6\linewidth}}{
        \cellcolor[rgb]{0.9,0.9,0.9}{
            \makecell[{{p{\linewidth}}}]{
                \texttt{\tiny{[P1$\langle$GM]}}
                \texttt{OUTPUT: (({[}2, 1{]}, 5), True)} \\
            }
        }
    }
     \\ \\

    \theutterance \stepcounter{utterance}  
    & \multicolumn{4}{p{0.6\linewidth}}{
        \cellcolor[rgb]{0.561,0.737,0.561}{
            \makecell[{{p{\linewidth}}}]{
                \texttt{\tiny{[P1$\rangle$GM]}}
                \texttt{\{} \\
\texttt{  "current\_mode": "probing",} \\
\texttt{  "input\_rationale": "Let's analyze the examples again. ({[}5, {-}5, 5{]}, 50) {-}> True. ({[}1, 2, 3{]}, 14) {-}> True. ({[}2, 1{]}, 5) {-}> True. ({[}1, 1{]}, 3) {-}> False. ({[}{-}1, 0, 1{]}, 3) {-}> False. ({[}0{]}, 0) {-}> True. ({[}{]}, {-}1) {-}> False. It seems like the function checks if the target can be formed by some combination of the numbers using addition and subtraction? Let's check: 5 + ({-}5) + 5 = 5 != 50. 5 * ({-}5) * 5 = {-}125. Maybe it's about the sum of absolute values? 5+5+5=15 != 50. Maybe it's about the product of absolute values? 5*5*5=125. Maybe it's about concatenating digits? No. Let's look at ({[}1, 2, 3{]}, 14). 1+2+3=6. 1*2*3=6. 12+3=15? 1+23=24? Maybe it's about forming the number by concatenating the digits of the numbers in the list} \\
            }
        }
    }
    & & \\ \\

    \theutterance \stepcounter{utterance}  
    & & & \multicolumn{4}{p{0.6\linewidth}}{
        \cellcolor[rgb]{0.9,0.9,0.9}{
            \makecell[{{p{\linewidth}}}]{
                \texttt{\tiny{[P1$\langle$GM]}}
                \texttt{Game Violation: ParseError: Response is not valid JSON.} \\
\texttt{Output ONLY:} \\
\texttt{One JSON with confidence\_score, current\_mode, input\_rationale, and hypothesis keys ...} \\
\texttt{OR} \\
\texttt{SOLVE: ```python ... ```} \\
            }
        }
    }
     \\ \\

    \theutterance \stepcounter{utterance}  
    & \multicolumn{4}{p{0.6\linewidth}}{
        \cellcolor[rgb]{0.561,0.737,0.561}{
            \makecell[{{p{\linewidth}}}]{
                \texttt{\tiny{[P1$\rangle$GM]}}
                \texttt{\{} \\
\texttt{  "current\_mode": "probing",} \\
\texttt{  "input\_rationale": "I need to find a pattern that explains: ({[}5, {-}5, 5{]}, 50) {-}> True, ({[}1, 2, 3{]}, 14) {-}> True, ({[}2, 1{]}, 5) {-}> True, ({[}1, 1{]}, 3) {-}> False, ({[}{-}1, 0, 1{]}, 3) {-}> False, ({[}0{]}, 0) {-}> True, ({[}{]}, {-}1) {-}> False. Let's look at ({[}2, 1{]}, 5) {-}> True. 2+1=3, 2*1=2. 2\^{}1=2, 1\^{}2=1. Maybe it's about the sum of squares? 2\^{}2+1\^{}2=5. Yes! Let's check ({[}1, 2, 3{]}, 14). 1\^{}2+2\^{}2+3\^{}2 = 1+4+9=14. Yes! Let's check ({[}5, {-}5, 5{]}, 50). 5\^{}2+({-}5)\^{}2+5\^{}2 = 25+25+25=75 != 50. Wait. Maybe it's not sum of squares. Let's re{-}examine ({[}5, {-}5, 5{]}, 50).} \\
            }
        }
    }
    & & \\ \\

    \theutterance \stepcounter{utterance}  
    & & & \multicolumn{4}{p{0.6\linewidth}}{
        \cellcolor[rgb]{0.9,0.9,0.9}{
            \makecell[{{p{\linewidth}}}]{
                \texttt{\tiny{[P1$\langle$GM]}}
                \texttt{Game Violation: ParseError: Response is not valid JSON.} \\
\texttt{Output ONLY:} \\
\texttt{One JSON with confidence\_score, current\_mode, input\_rationale, and hypothesis keys ...} \\
\texttt{OR} \\
\texttt{SOLVE: ```python ... ```} \\
            }
        }
    }
     \\ \\

    \theutterance \stepcounter{utterance}  
    & \multicolumn{4}{p{0.6\linewidth}}{
        \cellcolor[rgb]{0.561,0.737,0.561}{
            \makecell[{{p{\linewidth}}}]{
                \texttt{\tiny{[P1$\rangle$GM]}}
                \texttt{\{} \\
\texttt{  "current\_mode": "probing",} \\
\texttt{  "input\_rationale": "I have a strong hypothesis: the function checks if the sum of the squares of the items equals the target. Let's verify: ({[}1, 2, 3{]}, 14) {-}> 1\^{}2+2\^{}2+3\^{}2 = 1+4+9=14 (True). ({[}2, 1{]}, 5) {-}> 2\^{}2+1\^{}2 = 4+1=5 (True). ({[}0{]}, 0) {-}> 0\^{}2=0 (True). ({[}{]}, {-}1) {-}> sum is 0, 0 != {-}1 (False). ({[}{-}1, 0, 1{]}, 3) {-}> ({-}1)\^{}2+0\^{}2+1\^{}2 = 1+0+1=2 != 3 (False). ({[}1, 1{]}, 3) {-}> 1\^{}2+1\^{}2=2 != 3 (False). The only remaining example is ({[}5, {-}5, 5{]}, 50). 5\^{}2+({-}5)\^{}2+5\^{}2 = 25+25+25=75. This does not equal 50. So the hypothesis is incorrect. Let me re{-}read the first example. Maybe it's not sum of squares. What if it's sum of absolute values? 5} \\
            }
        }
    }
    & & \\ \\

    \theutterance \stepcounter{utterance}  
    & & & \multicolumn{4}{p{0.6\linewidth}}{
        \cellcolor[rgb]{0.9,0.9,0.9}{
            \makecell[{{p{\linewidth}}}]{
                \texttt{\tiny{[P1$\langle$GM]}}
                \texttt{Game Violation: ParseError: Response is not valid JSON.} \\
\texttt{Output ONLY:} \\
\texttt{One JSON with confidence\_score, current\_mode, input\_rationale, and hypothesis keys ...} \\
\texttt{OR} \\
\texttt{SOLVE: ```python ... ```} \\
            }
        }
    }
     \\ \\

    \theutterance \stepcounter{utterance}  
    & \multicolumn{4}{p{0.6\linewidth}}{
        \cellcolor[rgb]{0.561,0.737,0.561}{
            \makecell[{{p{\linewidth}}}]{
                \texttt{\tiny{[P1$\rangle$GM]}}
                \texttt{\{} \\
\texttt{  "current\_mode": "probing",} \\
\texttt{  "input\_rationale": "The sum of squares hypothesis failed for ({[}5, {-}5, 5{]}, 50) because 25+25+25=75 != 50. Let's look for another pattern. Maybe it's about the sum of the items multiplied by something? Or maybe it's about the product? Let's try a new probe to distinguish between sum of squares and other possibilities. Let's try ({[}2, 2{]}, 8). If it's sum of squares, 4+4=8 (True). If it's sum, 2+2=4 (False). If it's product, 2*2=4 (False).",} \\
\texttt{  "hypothesis": "unknown",} \\
\texttt{  "confidence\_score": 40} \\
\texttt{\}} \\
            }
        }
    }
    & & \\ \\

    \theutterance \stepcounter{utterance}  
    & & & \multicolumn{4}{p{0.6\linewidth}}{
        \cellcolor[rgb]{0.9,0.9,0.9}{
            \makecell[{{p{\linewidth}}}]{
                \texttt{\tiny{[P1$\langle$GM]}}
                \texttt{OUTPUT: (({[}3, 2, 1{]}, 15), False)} \\
            }
        }
    }
     \\ \\

    \theutterance \stepcounter{utterance}  
    & \multicolumn{4}{p{0.6\linewidth}}{
        \cellcolor[rgb]{0.561,0.737,0.561}{
            \makecell[{{p{\linewidth}}}]{
                \texttt{\tiny{[P1$\rangle$GM]}}
                \texttt{\{} \\
\texttt{  "current\_mode": "probing",} \\
\texttt{  "input\_rationale": "The previous hypothesis (sum of squares) predicted ({[}2, 2{]}, 8) {-}> True, but I haven't received that output yet. I received ({[}3, 2, 1{]}, 15) {-}> False. Sum of squares: 9+4+1=14 != 15. Sum: 3+2+1=6 != 15. Product: 6 != 15. This doesn't help much yet. Let's go back to the sum of squares idea. It worked for ({[}1, 2, 3{]}, 14) and ({[}2, 1{]}, 5). It failed for ({[}5, {-}5, 5{]}, 50). Maybe the first example is a trick or I'm misinterpreting it. Let's test ({[}2, 2{]}, 8) to confirm sum of squares.",} \\
\texttt{  "hypothesis": "unknown",} \\
\texttt{  "confidence\_score": 40} \\
\texttt{\}} \\
            }
        }
    }
    & & \\ \\

    \theutterance \stepcounter{utterance}  
    & & & \multicolumn{4}{p{0.6\linewidth}}{
        \cellcolor[rgb]{0.9,0.9,0.9}{
            \makecell[{{p{\linewidth}}}]{
                \texttt{\tiny{[P1$\langle$GM]}}
                \texttt{OUTPUT: (({[}10, {-}1, {-}1, 10{]}, 200), True)} \\
            }
        }
    }
     \\ \\

    \theutterance \stepcounter{utterance}  
    & \multicolumn{4}{p{0.6\linewidth}}{
        \cellcolor[rgb]{0.561,0.737,0.561}{
            \makecell[{{p{\linewidth}}}]{
                \texttt{\tiny{[P1$\rangle$GM]}}
                \texttt{\{} \\
\texttt{  "current\_mode": "probing",} \\
\texttt{  "input\_rationale": "Let's check the sum of squares hypothesis again with the new example: ({[}10, {-}1, {-}1, 10{]}, 200). Sum of squares: 10\^{}2 + ({-}1)\^{}2 + ({-}1)\^{}2 + 10\^{}2 = 100 + 1 + 1 + 100 = 202. The output is True, but 202 != 200. So sum of squares is incorrect. Let's look at the numbers: 10, {-}1, {-}1, 10. Target 200. 10 * 10 * ({-}1) * ({-}1) = 100. 10 * 10 + ({-}1) * ({-}1) = 101. 10 * 10 * 2 = 200? Where does 2 come from? Maybe it's the count of positive numbers? No. Let's look at ({[}5, {-}5, 5{]}, 50). 5 * 5 * 2 = 50? Count of positive numbers is 2. 5*5*2=50. For ({[}10, {-}1, {-}1, 10{]}, 200): Positive numbers are 10,} \\
            }
        }
    }
    & & \\ \\

    \theutterance \stepcounter{utterance}  
    & & & \multicolumn{4}{p{0.6\linewidth}}{
        \cellcolor[rgb]{0.9,0.9,0.9}{
            \makecell[{{p{\linewidth}}}]{
                \texttt{\tiny{[P1$\langle$GM]}}
                \texttt{Game Violation: ParseError: Response is not valid JSON.} \\
\texttt{Output ONLY:} \\
\texttt{One JSON with confidence\_score, current\_mode, input\_rationale, and hypothesis keys ...} \\
\texttt{OR} \\
\texttt{SOLVE: ```python ... ```} \\
            }
        }
    }
     \\ \\

    \theutterance \stepcounter{utterance}  
    & \multicolumn{4}{p{0.6\linewidth}}{
        \cellcolor[rgb]{0.561,0.737,0.561}{
            \makecell[{{p{\linewidth}}}]{
                \texttt{\tiny{[P1$\rangle$GM]}}
                \texttt{\{} \\
\texttt{  "current\_mode": "probing",} \\
\texttt{  "input\_rationale": "I need to find a pattern that fits all examples. Let's list them: ({[}5, {-}5, 5{]}, 50) {-}> True, ({[}{]}, {-}1) {-}> False, ({[}0{]}, 0) {-}> True, ({[}1, 1{]}, 3) {-}> False, ({[}1, 2, 3{]}, 14) {-}> True, ({[}{-}1, 0, 1{]}, 3) {-}> False, ({[}2, 1{]}, 5) {-}> True, ({[}3, 2, 1{]}, 15) {-}> False, ({[}10, {-}1, {-}1, 10{]}, 200) {-}> True. Let's check sum of squares again: ({[}1,2,3{]} {-}> 14 True), ({[}2,1{]} {-}> 5 True), ({[}0{]} {-}> 0 True), ({[}1,1{]} {-}> 2 != 3 False), ({[}{-}1,0,1{]} {-}> 2 != 3 False), ({[}3,2,1{]} {-}> 14 != 15 False). This fits all except ({[}5,{-}5,5{]} {-}> 75 != 50) and ({[}10,{-}1,{-}1,10{]} {-}> 202 != 200). Wait, 75 is close to 5} \\
            }
        }
    }
    & & \\ \\

    \theutterance \stepcounter{utterance}  
    & & & \multicolumn{4}{p{0.6\linewidth}}{
        \cellcolor[rgb]{0.9,0.9,0.9}{
            \makecell[{{p{\linewidth}}}]{
                \texttt{\tiny{[P1$\langle$GM]}}
                \texttt{Game Violation: ParseError: Response is not valid JSON.} \\
\texttt{Output ONLY:} \\
\texttt{One JSON with confidence\_score, current\_mode, input\_rationale, and hypothesis keys ...} \\
\texttt{OR} \\
\texttt{SOLVE: ```python ... ```} \\
            }
        }
    }
     \\ \\

    \theutterance \stepcounter{utterance}  
    & \multicolumn{4}{p{0.6\linewidth}}{
        \cellcolor[rgb]{0.561,0.737,0.561}{
            \makecell[{{p{\linewidth}}}]{
                \texttt{\tiny{[P1$\rangle$GM]}}
                \texttt{\{} \\
\texttt{  "current\_mode": "probing",} \\
\texttt{  "input\_rationale": "I have a strong hypothesis: the function checks if the sum of the squares of the items equals the target. Let's verify: ({[}1, 2, 3{]}, 14) {-}> 1\^{}2+2\^{}2+3\^{}2 = 1+4+9=14 (True). ({[}2, 1{]}, 5) {-}> 2\^{}2+1\^{}2 = 4+1=5 (True). ({[}0{]}, 0) {-}> 0\^{}2=0 (True). ({[}{]}, {-}1) {-}> sum is 0, 0 != {-}1 (False). ({[}{-}1, 0, 1{]}, 3) {-}> ({-}1)\^{}2+0\^{}2+1\^{}2 = 1+0+1=2 != 3 (False). ({[}1, 1{]}, 3) {-}> 1\^{}2+1\^{}2=2 != 3 (False). ({[}3, 2, 1{]}, 15) {-}> 3\^{}2+2\^{}2+1\^{}2 = 9+4+1=14 != 15 (False). The only remaining examples are ({[}5, {-}5, 5{]}, 50) and ({[}10, {-}1, {-}1, 10{]}, 200). For ({[}5, {-}} \\
            }
        }
    }
    & & \\ \\

    \theutterance \stepcounter{utterance}  
    & & & \multicolumn{2}{p{0.3\linewidth}}{
        \cellcolor[rgb]{0.9,0.9,0.9}{
            \makecell[{{p{\linewidth}}}]{
                \texttt{\tiny{[GM$|$GM]}}
                \texttt{Game Verdict: LOSS (turn limit reached)} \\
            }
        }
    }
    & & \\ \\

\end{supertabular}
}

\end{document}
